% Options for packages loaded elsewhere
\PassOptionsToPackage{unicode}{hyperref}
\PassOptionsToPackage{hyphens}{url}
\documentclass[
  ukrainian,
]{article}
\usepackage{xcolor}
\usepackage[margin=2.5cm]{geometry}
\usepackage{amsmath,amssymb}
\setcounter{secnumdepth}{-\maxdimen} % remove section numbering
\usepackage{iftex}
\ifPDFTeX
  \usepackage[T1]{fontenc}
  \usepackage[utf8]{inputenc}
  \usepackage{textcomp} % provide euro and other symbols
\else % if luatex or xetex
  \usepackage{unicode-math} % this also loads fontspec
  \defaultfontfeatures{Scale=MatchLowercase}
  \defaultfontfeatures[\rmfamily]{Ligatures=TeX,Scale=1}
\fi
\usepackage{lmodern}
\ifPDFTeX\else
\setmainfont{DejaVu Serif}
\setsansfont{DejaVu Sans}
\setmonofont{DejaVu Sans Mono}
\setmathfont{STIX Two Math}
\fi
\ifPDFTeX\else
  % xetex/luatex font selection
\fi
% Use upquote if available, for straight quotes in verbatim environments
\IfFileExists{upquote.sty}{\usepackage{upquote}}{}
\IfFileExists{microtype.sty}{% use microtype if available
  \usepackage[]{microtype}
  \UseMicrotypeSet[protrusion]{basicmath} % disable protrusion for tt fonts
}{}
\makeatletter
\@ifundefined{KOMAClassName}{% if non-KOMA class
  \IfFileExists{parskip.sty}{%
    \usepackage{parskip}
  }{% else
    \setlength{\parindent}{0pt}
    \setlength{\parskip}{6pt plus 2pt minus 1pt}}
}{% if KOMA class
  \KOMAoptions{parskip=half}}
\makeatother
\usepackage{graphicx}
\graphicspath{{/Users/kalinicenkovitalijmikolajovic/.codex/worktrees/5c15/Script_for_translate_book/Output_Final/}}
\setkeys{Gin}{width=0.74\linewidth,height=0.78\textheight,keepaspectratio}
\makeatletter
\newsavebox\pandoc@box
\newcommand*\pandocbounded[1]{% scales image to fit in text height/width
  \sbox\pandoc@box{#1}%
  \Gscale@div\@tempa{\textheight}{\dimexpr\ht\pandoc@box+\dp\pandoc@box\relax}%
  \Gscale@div\@tempb{\linewidth}{\wd\pandoc@box}%
  \ifdim\@tempb\p@<\@tempa\p@\let\@tempa\@tempb\fi% select the smaller of both
  \ifdim\@tempa\p@<\p@\scalebox{\@tempa}{\usebox\pandoc@box}%
  \else\usebox{\pandoc@box}%
  \fi%
}
% Set default figure placement to htbp
\def\fps@figure{htbp}
\makeatother
\ifLuaTeX
\usepackage[bidi=basic,shorthands=off,provide=*]{babel}
\else
\usepackage[bidi=default,shorthands=off,provide=*]{babel}
\fi
\ifLuaTeX
  \usepackage{selnolig} % disable illegal ligatures
\fi
\setlength{\emergencystretch}{3em} % prevent overfull lines
\providecommand{\tightlist}{%
  \setlength{\itemsep}{0pt}\setlength{\parskip}{0pt}}
\usepackage{bookmark}
\IfFileExists{xurl.sty}{\usepackage{xurl}}{} % add URL line breaks if available
\urlstyle{same}
\hypersetup{
  pdflang={uk},
  hidelinks,
  pdfcreator={LaTeX via pandoc}}

\author{}
\date{}

\begin{document}

\subsection{Лекція 47}\label{ux43bux435ux43aux446ux456ux44f-47}

\section{РОЗВ'ЯЗАННЯ ДВОВИМІРНИХ ЗАДАЧ МЕТОДОМ КОНФОРМНИХ
ВІДОБРАЖЕНЬ}\label{ux440ux43eux437ux432ux44fux437ux430ux43dux43dux44f-ux434ux432ux43eux432ux438ux43cux456ux440ux43dux438ux445-ux437ux430ux434ux430ux447-ux43cux435ux442ux43eux434ux43eux43c-ux43aux43eux43dux444ux43eux440ux43cux43dux438ux445-ux432ux456ux434ux43eux431ux440ux430ux436ux435ux43dux44c}

МЕТА ЛЕКЦІЇ: Показати, як деякі двовимірні крайові задачі для рівняння
Лапласа можна звести до простіших задач за допомогою конформного
відображення.

\pandocbounded{\includegraphics[keepaspectratio]{images/_page_353_Figure_9.jpeg}}

Ідея методу полягає в тому, що складну область на площині \((x,y)\)
замінюють простішою областю на площині \((u,v)\) за допомогою
комплексного відображення \(w=f(z)\), де

\[
z = x + iy, \qquad w = u + iv.
\]

Для конформних відображень гармонічність зберігається: якщо
\(\varphi(x,y)\) задовольняє рівняння Лапласа в початковій області, то
після переходу до нових координат ми знову дістаємо рівняння Лапласа,
але вже в геометрично простішій області.

У цій лекції розглянуто:

\begin{enumerate}
\def\labelenumi{\arabic{enumi}.}
\tightlist
\item
  основні властивості конформних відображень;
\item
  приклади побудови таких відображень;
\item
  застосування методу до задач Діріхле.
\end{enumerate}

\section{Комплексні функції та конформні
відображення}\label{ux43aux43eux43cux43fux43bux435ux43aux441ux43dux456-ux444ux443ux43dux43aux446ux456ux457-ux442ux430-ux43aux43eux43dux444ux43eux440ux43cux43dux456-ux432ux456ux434ux43eux431ux440ux430ux436ux435ux43dux43dux44f}

Комплексну функцію записують у вигляді

\[
w = f(z).
\]

Вона переводить точки площини \(z\) у точки площини \(w\). Наприклад,
відображення

\[
w = z^2
\]

переносить перший квадрант площини \(z\) у верхню півплощину площини
\(w\).

\pandocbounded{\includegraphics[keepaspectratio]{images/_page_355_Figure_5.jpeg}}

Рис. 47.1. Площина комплексної змінної \(z\).

\pandocbounded{\includegraphics[keepaspectratio]{images/_page_355_Figure_7.jpeg}}

Рис. 47.2. Відображення першого квадранта площини \(z\) у верхню
півплощину площини \(w\).

Щоб побачити це в дійсних координатах, підставимо \(z=x+iy\) у формулу
\(w=z^2\):

\[
w = (x+iy)^2 = x^2 - y^2 + 2ixy.
\]

Отже,

\[
u = x^2 - y^2, \qquad v = 2xy.
\]

Це і є дійсна форма того самого відображення. Вона показує, як криві у
площині \((x,y)\) переходять у криві в площині \((u,v)\).

\section{Означення конформного
відображення}\label{ux43eux437ux43dux430ux447ux435ux43dux43dux44f-ux43aux43eux43dux444ux43eux440ux43cux43dux43eux433ux43e-ux432ux456ux434ux43eux431ux440ux430ux436ux435ux43dux43dux44f}

Відображення \(w=f(z)\) називається конформним у точці \(z_0\), якщо
\(f'(z_0)\neq 0\). Воно називається конформним в області \(D\), якщо ця
умова виконується в кожній точці області.

Геометрично це означає, що відображення зберігає кути між кривими. Саме
ця властивість робить його корисним для розв'язування крайових задач:
складну межу можна перенести в простішу, не зруйнувавши структуру
задачі.

\section{Рівняння Лапласа у верхній
півплощині}\label{ux440ux456ux432ux43dux44fux43dux43dux44f-ux43bux430ux43fux43bux430ux441ux430-ux443-ux432ux435ux440ux445ux43dux456ux439-ux43fux456ux432ux43fux43bux43eux449ux438ux43dux456}

Розглянемо задачу Діріхле у верхній півплощині. Основна ідея полягає в
тому, щоб підібрати відображення, яке переведе верхню півплощину у
горизонтальну смугу. Для цього використовується класичне перетворення

\[
w = \ln\left(\frac{z-1}{z+1}\right).
\]

Це відображення переносить верхню півплощину у смугу

\[
-\infty < u < \infty, \qquad 0 < v < \pi.
\]

\pandocbounded{\includegraphics[keepaspectratio]{images/_page_357_Figure_11.jpeg}}

Рис. 47.3. Крайова задача у верхній півплощині.

\pandocbounded{\includegraphics[keepaspectratio]{images/_page_358_Figure_2.jpeg}}

Рис. 47.4. Перетворення складної задачі на простішу за допомогою
конформного відображення.

У смузі розв'язок уже знаходиться елементарно. Якщо на одній межі задано
нульовий потенціал, а на іншій одиничний, то

\[
\varphi(u,v) = \frac{v}{\pi}.
\]

Далі потрібно виразити \(v\) через початкові координати \(x\) і \(y\) та
підставити цей вираз у знайдену формулу.

\section{Область між двома
колами}\label{ux43eux431ux43bux430ux441ux442ux44c-ux43cux456ux436-ux434ux432ux43eux43cux430-ux43aux43eux43bux430ux43cux438}

Другий типовий приклад полягає у відображенні області між двома
неконцентричними колами на кільце. Після такого перетворення задача
Діріхле зводиться до стандартної задачі в кільці, яку вже можна
розв'язувати у полярних координатах.

\pandocbounded{\includegraphics[keepaspectratio]{images/_page_359_Figure_10.jpeg}}

Рис. 47.5. Конформне відображення області на кільце.

Суть методу лишається тією самою:

\begin{enumerate}
\def\labelenumi{\arabic{enumi}.}
\tightlist
\item
  знайти зручне конформне відображення;
\item
  перенести задачу у простішу область;
\item
  розв'язати її в нових координатах;
\item
  повернутися до початкових змінних.
\end{enumerate}

Саме тому конформні відображення є одним із найсильніших інструментів
для двовимірних гармонічних задач.

\end{document}
